% Options for packages loaded elsewhere
\PassOptionsToPackage{unicode}{hyperref}
\PassOptionsToPackage{hyphens}{url}
\documentclass[
  ignorenonframetext,
]{beamer}
\newif\ifbibliography
\usepackage{pgfpages}
\setbeamertemplate{caption}[numbered]
\setbeamertemplate{caption label separator}{: }
\setbeamercolor{caption name}{fg=normal text.fg}
\beamertemplatenavigationsymbolsempty
% remove section numbering
\setbeamertemplate{part page}{
  \centering
  \begin{beamercolorbox}[sep=16pt,center]{part title}
    \usebeamerfont{part title}\insertpart\par
  \end{beamercolorbox}
}
\setbeamertemplate{section page}{
  \centering
  \begin{beamercolorbox}[sep=12pt,center]{section title}
    \usebeamerfont{section title}\insertsection\par
  \end{beamercolorbox}
}
\setbeamertemplate{subsection page}{
  \centering
  \begin{beamercolorbox}[sep=8pt,center]{subsection title}
    \usebeamerfont{subsection title}\insertsubsection\par
  \end{beamercolorbox}
}
% Prevent slide breaks in the middle of a paragraph
\widowpenalties 1 10000
\raggedbottom
\AtBeginPart{
  \frame{\partpage}
}
\AtBeginSection{
  \ifbibliography
  \else
    \frame{\sectionpage}
  \fi
}
\AtBeginSubsection{
  \frame{\subsectionpage}
}
\usepackage{iftex}
\ifPDFTeX
  \usepackage[T1]{fontenc}
  \usepackage[utf8]{inputenc}
  \usepackage{textcomp} % provide euro and other symbols
\else % if luatex or xetex
  \usepackage{unicode-math} % this also loads fontspec
  \defaultfontfeatures{Scale=MatchLowercase}
  \defaultfontfeatures[\rmfamily]{Ligatures=TeX,Scale=1}
\fi
\usepackage{lmodern}
\ifPDFTeX\else
  % xetex/luatex font selection
\fi
% Use upquote if available, for straight quotes in verbatim environments
\IfFileExists{upquote.sty}{\usepackage{upquote}}{}
\IfFileExists{microtype.sty}{% use microtype if available
  \usepackage[]{microtype}
  \UseMicrotypeSet[protrusion]{basicmath} % disable protrusion for tt fonts
}{}
\makeatletter
\@ifundefined{KOMAClassName}{% if non-KOMA class
  \IfFileExists{parskip.sty}{%
    \usepackage{parskip}
  }{% else
    \setlength{\parindent}{0pt}
    \setlength{\parskip}{6pt plus 2pt minus 1pt}}
}{% if KOMA class
  \KOMAoptions{parskip=half}}
\makeatother
\usepackage{color}
\usepackage{fancyvrb}
\newcommand{\VerbBar}{|}
\newcommand{\VERB}{\Verb[commandchars=\\\{\}]}
\DefineVerbatimEnvironment{Highlighting}{Verbatim}{commandchars=\\\{\}}
% Add ',fontsize=\small' for more characters per line
\newenvironment{Shaded}{}{}
\newcommand{\AlertTok}[1]{\textcolor[rgb]{1.00,0.00,0.00}{\textbf{#1}}}
\newcommand{\AnnotationTok}[1]{\textcolor[rgb]{0.38,0.63,0.69}{\textbf{\textit{#1}}}}
\newcommand{\AttributeTok}[1]{\textcolor[rgb]{0.49,0.56,0.16}{#1}}
\newcommand{\BaseNTok}[1]{\textcolor[rgb]{0.25,0.63,0.44}{#1}}
\newcommand{\BuiltInTok}[1]{\textcolor[rgb]{0.00,0.50,0.00}{#1}}
\newcommand{\CharTok}[1]{\textcolor[rgb]{0.25,0.44,0.63}{#1}}
\newcommand{\CommentTok}[1]{\textcolor[rgb]{0.38,0.63,0.69}{\textit{#1}}}
\newcommand{\CommentVarTok}[1]{\textcolor[rgb]{0.38,0.63,0.69}{\textbf{\textit{#1}}}}
\newcommand{\ConstantTok}[1]{\textcolor[rgb]{0.53,0.00,0.00}{#1}}
\newcommand{\ControlFlowTok}[1]{\textcolor[rgb]{0.00,0.44,0.13}{\textbf{#1}}}
\newcommand{\DataTypeTok}[1]{\textcolor[rgb]{0.56,0.13,0.00}{#1}}
\newcommand{\DecValTok}[1]{\textcolor[rgb]{0.25,0.63,0.44}{#1}}
\newcommand{\DocumentationTok}[1]{\textcolor[rgb]{0.73,0.13,0.13}{\textit{#1}}}
\newcommand{\ErrorTok}[1]{\textcolor[rgb]{1.00,0.00,0.00}{\textbf{#1}}}
\newcommand{\ExtensionTok}[1]{#1}
\newcommand{\FloatTok}[1]{\textcolor[rgb]{0.25,0.63,0.44}{#1}}
\newcommand{\FunctionTok}[1]{\textcolor[rgb]{0.02,0.16,0.49}{#1}}
\newcommand{\ImportTok}[1]{\textcolor[rgb]{0.00,0.50,0.00}{\textbf{#1}}}
\newcommand{\InformationTok}[1]{\textcolor[rgb]{0.38,0.63,0.69}{\textbf{\textit{#1}}}}
\newcommand{\KeywordTok}[1]{\textcolor[rgb]{0.00,0.44,0.13}{\textbf{#1}}}
\newcommand{\NormalTok}[1]{#1}
\newcommand{\OperatorTok}[1]{\textcolor[rgb]{0.40,0.40,0.40}{#1}}
\newcommand{\OtherTok}[1]{\textcolor[rgb]{0.00,0.44,0.13}{#1}}
\newcommand{\PreprocessorTok}[1]{\textcolor[rgb]{0.74,0.48,0.00}{#1}}
\newcommand{\RegionMarkerTok}[1]{#1}
\newcommand{\SpecialCharTok}[1]{\textcolor[rgb]{0.25,0.44,0.63}{#1}}
\newcommand{\SpecialStringTok}[1]{\textcolor[rgb]{0.73,0.40,0.53}{#1}}
\newcommand{\StringTok}[1]{\textcolor[rgb]{0.25,0.44,0.63}{#1}}
\newcommand{\VariableTok}[1]{\textcolor[rgb]{0.10,0.09,0.49}{#1}}
\newcommand{\VerbatimStringTok}[1]{\textcolor[rgb]{0.25,0.44,0.63}{#1}}
\newcommand{\WarningTok}[1]{\textcolor[rgb]{0.38,0.63,0.69}{\textbf{\textit{#1}}}}
\setlength{\emergencystretch}{3em} % prevent overfull lines
\providecommand{\tightlist}{%
  \setlength{\itemsep}{0pt}\setlength{\parskip}{0pt}}
\usepackage{bookmark}
\IfFileExists{xurl.sty}{\usepackage{xurl}}{} % add URL line breaks if available
\urlstyle{same}
\hypersetup{
  hidelinks,
  pdfcreator={LaTeX via pandoc}}

\author{\texorpdfstring{}{}}
\date{}

\begin{document}

\begin{frame}[fragile]{Appendix: Mathematical and Algorithmic Details
\label{sec:technical_appendix}}
\protect\phantomsection\label{appendix-mathematical-and-algorithmic-details}
\emph{This appendix collects the formal mathematical definitions,
derivations, and algorithmic specifications referenced from the main
methodology section.}

\begin{block}{Citation-Weighted Hypothesis Scoring Formula
\label{sec:appendix_scoring}}
\protect\phantomsection\label{citation-weighted-hypothesis-scoring-formula}
For each hypothesis \(H\), we compute a citation-weighted evidence score
aggregating all assertions relevant to \(H\):

\begin{equation}
\text{score}(H) = \frac{\sum_{a \in S(H)} w(a) - \sum_{a \in C(H)} w(a)}{\sum_{a \in A(H)} w(a)} \label{eq:app_score}
\end{equation}

where \(S(H)\) is the set of supporting assertions, \(C(H)\) is the set
of contradicting assertions, \(A(H)\) is all assertions for \(H\)
(including neutral), and the weight function is:

\begin{equation}
w(a) = \log(1 + \text{citations}(a)) \cdot \text{confidence}(a) \label{eq:app_weight}
\end{equation}

The logarithmic citation weighting ensures that highly cited papers
carry more influence while preventing any single blockbuster paper from
dominating the score. The score lies in \([-1, 1]\): values near \(+1\)
indicate strong supporting evidence, values near \(-1\) indicate strong
contradicting evidence, and values near \(0\) indicate balanced or
insufficient evidence. Crucially, as discussed in the main text, this
score represents a \emph{relative evidentiary ranking} within the
current literature topology, not a calibrated Bayesian probability of
the hypothesis actually being true.

\textbf{Temporal aggregation.} We additionally compute temporal trends
by evaluating the cumulative score at each year \(t\), using only
assertions from papers published in year \(\leq t\):

\begin{equation}
\text{score}(H, t) = \frac{\sum_{a \in S(H,t)} w(a) - \sum_{a \in C(H,t)} w(a)}{\sum_{a \in A(H,t)} w(a)} \label{eq:app_score_t}
\end{equation}

This reveals whether support for a hypothesis is growing, declining, or
plateauing over time.
\end{block}

\begin{block}{Non-negative Matrix Factorization (NMF) for Topic Modeling
\label{sec:appendix_nmf}}
\protect\phantomsection\label{non-negative-matrix-factorization-nmf-for-topic-modeling}
We apply NMF to the TF-IDF matrix of the corpus to discover latent
topics. Given the document-term matrix
\(V \in \mathbb{R}^{n \times m}_{\geq 0}\), NMF finds factor matrices
\(W \in \mathbb{R}^{n \times k}_{\geq 0}\) and
\(H \in \mathbb{R}^{k \times m}_{\geq 0}\) such that \(V \approx WH\),
where \(k\) is the number of topics.

We use multiplicative update rules \citep{lee1999nmf}:

\begin{equation}
H \leftarrow H \odot \frac{W^T V}{W^T W H + \epsilon}, \quad W \leftarrow W \odot \frac{V H^T}{W H H^T + \epsilon} \label{eq:nmf_update}
\end{equation}

with \(\epsilon = 10^{-10}\) for numerical stability and a fixed random
seed of 42 for reproducibility (ensuring deterministic topic alignment
across pipeline runs, with empirical stability confirmed via Jaccard
similarities \(> 0.90\) across alternative seeds).

\textbf{Term-Frequency Inverse Document Frequency (TF-IDF).} The
document-term matrix is constructed using TF-IDF weighting
\citep{salton1975vector}. For term \(t\) in document \(d\):

\begin{equation}
\text{TF-IDF}(t, d) = \text{tf}(t, d) \cdot \log\!\left(\frac{N}{\text{df}(t)}\right) \label{eq:tfidf}
\end{equation}

where \(\text{tf}(t, d)\) is the term frequency, \(N\) is the total
number of documents, and \(\text{df}(t)\) is the document frequency of
term \(t\).
\end{block}

\begin{block}{Field Growth-Rate Estimation \label{sec:appendix_growth}}
\protect\phantomsection\label{field-growth-rate-estimation}
The \textbf{mean year-over-year growth rate} \(\bar{g}\) is the
arithmetic mean of annual growth rates computed only for years where the
prior year had non-zero publications:

\begin{equation}
\bar{g} = \frac{1}{|Y|} \sum_{y \in Y} \frac{n_y - n_{y-1}}{n_{y-1}} \label{eq:mean_growth}
\end{equation}

where \(Y = \{y : n_{y-1} > 0\}\) and \(n_y\) is the number of
publications in year \(y\).

The \textbf{doubling time} \(t_d\) is derived from the mean annual
growth rate:

\begin{equation}
t_d = \frac{\ln 2}{\ln(1 + \bar{g})} \label{eq:doubling_time}
\end{equation}

The \textbf{compound annual growth rate} (CAGR) over the full span
\([y_0, y_T]\) is:

\begin{equation}
\text{CAGR} = \left(\frac{n_{\text{cumulative}}(y_T)}{n_{\text{cumulative}}(y_0)}\right)^{1/(y_T - y_0)} - 1 \label{eq:cagr}
\end{equation}

For the current corpus, CAGR \(= 16.99\%\). The more recent growth phase
(2010--2026) exhibits substantially higher annualized growth.
\end{block}

\begin{block}{Advanced Visualization Methods \label{sec:appendix_viz}}
\protect\phantomsection\label{advanced-visualization-methods}
\begin{block}{PCA of TF-IDF Embeddings}
\protect\phantomsection\label{pca-of-tf-idf-embeddings}
Principal Component Analysis (PCA) is applied to the TF-IDF matrix \(V\)
to project each document into a 2-D space. The projection preserves the
directions of maximum variance, enabling visual inspection of document
clustering by domain. Loading arrows overlay the top-variance terms onto
the scatter plot, showing which vocabulary drives the principal
components.
\end{block}

\begin{block}{Hierarchical Clustering Dendrogram}
\protect\phantomsection\label{hierarchical-clustering-dendrogram}
For each domain \(s\), we compute the centroid
\(\bar{v}_s = \frac{1}{|D_s|} \sum_{d \in D_s} v_d\) where \(D_s\) is
the set of documents in domain \(s\) and \(v_d\) is the TF-IDF vector of
document \(d\). Ward linkage is applied to the centroid matrix to
produce a hierarchical clustering dendrogram showing semantic proximity
between domains.
\end{block}

\begin{block}{Term Heatmap}
\protect\phantomsection\label{term-heatmap}
For each domain \(s\) and term \(t\), we compute the mean TF-IDF weight
\(\bar{w}_{s,t} = \frac{1}{|D_s|} \sum_{d \in D_s} \text{TF-IDF}(t, d)\).
The heatmap displays \(\bar{w}_{s,t}\) for the top-\(k\) terms (by
global document frequency) across all domains, with cell intensity
proportional to mean weight. This reveals distinctive vocabulary
patterns that differentiate domains beyond the keyword-level
classification used for subfield assignment.
\end{block}

\begin{block}{Term Co-occurrence Matrix}
\protect\phantomsection\label{term-co-occurrence-matrix}
The co-occurrence matrix \(C \in \mathbb{R}^{k \times k}\) counts the
number of documents in which two terms appear together. For top-\(k\)
terms by document frequency,
\(C_{ij} = |\{d : t_i \in d \land t_j \in d\}|\). The matrix is
normalized to \([0, 1]\) by dividing by the maximum entry and visualized
as a symmetric heatmap.
\end{block}
\end{block}

\begin{block}{Nanopublication RDF Schema \label{sec:appendix_rdf}}
\protect\phantomsection\label{nanopublication-rdf-schema}
Each nanopublication is serialized to RDF/TriG per the nanopublication
standard \citep{groth2010anatomy, kuhn2016decentralized}, producing four
named graphs. The following annotated example illustrates the structure
for a single assertion:

\begin{Shaded}
\begin{Highlighting}[]
\NormalTok{@prefix np: \textless{}http://www.nanopub.org/nschema\#\textgreater{} .}
\NormalTok{@prefix prov: \textless{}http://www.w3.org/ns/prov\#\textgreater{} .}
\NormalTok{@prefix dc: \textless{}http://purl.org/dc/terms/\textgreater{} .}
\NormalTok{@prefix aif: \textless{}http://activeinference.institute/ontology/\textgreater{} .}
\NormalTok{@prefix xsd: \textless{}http://www.w3.org/2001/XMLSchema\#\textgreater{} .}

\NormalTok{\# HEAD GRAPH: links nanopub to its three component graphs}
\NormalTok{\textless{}http://activeinference.institute/nanopub/a1b2c3d4e5f6\#head\textgreater{} \{}
\NormalTok{  \textless{}http://activeinference.institute/nanopub/a1b2c3d4e5f6\textgreater{}}
\NormalTok{    a np:Nanopublication ;}
\NormalTok{    np:hasAssertion   \textless{}...\#assertion\textgreater{} ;}
\NormalTok{    np:hasProvenance   \textless{}...\#provenance\textgreater{} ;}
\NormalTok{    np:hasPublicationInfo \textless{}...\#pubinfo\textgreater{} .}
\NormalTok{\}}

\NormalTok{\# ASSERTION GRAPH: the core scientific claim}
\NormalTok{\textless{}http://activeinference.institute/nanopub/a1b2c3d4e5f6\#assertion\textgreater{} \{}
\NormalTok{  aif:paper/10.1038\_nrn2787 aif:asserts aif:assertion/a1b2c3 .}
\NormalTok{  aif:assertion/a1b2c3}
\NormalTok{    aif:supports aif:hypothesis/fep\_universality ;}
\NormalTok{    aif:claim "The paper provides foundational support for FEP as a}
\NormalTok{               unified brain theory."\^{}\^{}xsd:string ;}
\NormalTok{    aif:confidence "0.85"\^{}\^{}xsd:double ;}
\NormalTok{    aif:citationCount "5000"\^{}\^{}xsd:integer .}
\NormalTok{\}}

\NormalTok{\# PROVENANCE GRAPH: extraction lineage}
\NormalTok{\textless{}http://activeinference.institute/nanopub/a1b2c3d4e5f6\#provenance\textgreater{} \{}
\NormalTok{  aif:assertion/a1b2c3}
\NormalTok{    prov:wasGeneratedBy  \textless{}http://activeinference.institute/nanopub/a1b2c3d4e5f6\textgreater{} ;}
\NormalTok{    prov:generatedAtTime "2026{-}01{-}15T12:00:00+00:00"\^{}\^{}xsd:dateTime ;}
\NormalTok{    prov:wasAttributedTo "act\_inf\_metaanalysis/gemma3:4b"\^{}\^{}xsd:string ;}
\NormalTok{    prov:hadPrimarySource aif:paper/10.1038\_nrn2787 .}
\NormalTok{\}}

\NormalTok{\# PUBLICATION INFO GRAPH: nanopublication metadata}
\NormalTok{\textless{}http://activeinference.institute/nanopub/a1b2c3d4e5f6\#pubinfo\textgreater{} \{}
\NormalTok{  \textless{}http://activeinference.institute/nanopub/a1b2c3d4e5f6\textgreater{}}
\NormalTok{    dc:created "2026{-}01{-}15T12:00:00+00:00"\^{}\^{}xsd:dateTime ;}
\NormalTok{    dc:creator "act\_inf\_metaanalysis/gemma3:4b"\^{}\^{}xsd:string ;}
\NormalTok{    dc:license \textless{}https://creativecommons.org/publicdomain/zero/1.0/\textgreater{} .}
\NormalTok{\}}
\end{Highlighting}
\end{Shaded}

\begin{block}{Namespace Definitions}
\protect\phantomsection\label{namespace-definitions}
\begin{table}[htbp]
\centering
\caption{RDF namespace definitions used in the knowledge graph and nanopublication serialization. Each prefix maps to a W3C or domain-specific URI.}
\label{tab:namespace_definitions}
\begin{tabular}{lll}
\toprule
\textbf{Prefix} & \textbf{URI} & \textbf{Purpose} \\
\midrule
\texttt{np:} & \texttt{http://www.nanopub.org/nschema\#} & Nanopub structural predicates \\
\texttt{prov:} & \texttt{http://www.w3.org/ns/prov\#} & PROV-O provenance model \\
\texttt{dc:} & \texttt{http://purl.org/dc/terms/} & Dublin Core metadata \\
\texttt{aif:} & \texttt{http://activeinference.institute/ontology/} & Domain-specific predicates \\
\texttt{xsd:} & \texttt{http://www.w3.org/2001/XMLSchema\#} & XML Schema datatypes \\
\bottomrule
\end{tabular}
\end{table}
\end{block}

\begin{block}{Core Triple Patterns}
\protect\phantomsection\label{core-triple-patterns}
The knowledge graph encodes five fundamental relationships:

\begin{table}[htbp]
\centering
\caption{Core RDF triple patterns encoding the five fundamental relationships in the knowledge graph. Each pattern links paper, assertion, hypothesis, or subfield nodes.}
\label{tab:core_triple_patterns}
\begin{tabular}{ll}
\toprule
\textbf{Triple Pattern} & \textbf{Meaning} \\
\midrule
\texttt{Paper --aif:asserts--> Assertion} & A paper makes a claim \\
\texttt{Paper --aif:cites--> Paper} & Intra-corpus citation link \\
\texttt{Paper --aif:belongsTo--> Subfield} & Domain classification \\
\texttt{Assertion --aif:supports--> Hypothesis} & Supporting evidence \\
\texttt{Assertion --aif:contradicts--> Hypothesis} & Contradicting evidence \\
\bottomrule
\end{tabular}
\end{table}
\end{block}
\end{block}
\end{frame}

\end{document}
